% lemmas_B1_B4.tex
% Agent B3, 2026-05-03.
% Closes Gap 1 (spatial kernel WF-set via Bessel J_0) and
% Gap 4 (IR Plancherel convergence near r=0) for the Bianchi VII_0
% SLE Hadamard construction of A4/note.tex.
%
% CITATION VERIFICATION LOG (arXiv API, 2026-05-03):
%   [AV13]  Avetisyan--Verch, arXiv:1212.6180, CQG 30 (2013) 155006. VERIFIED.
%   [BN23]  Banerjee--Niedermaier, arXiv:2305.11388, JMP 64 (2023) 113503. VERIFIED.
%   [TB13]  Them--Brum, arXiv:1302.3174, CQG 30 (2013) 235035. VERIFIED.
%   [Rad96] Radzikowski, CMP 179 (1996) 529-553. No arXiv; DOI 10.1007/BF02100096. VERIFIED.
%   [Ver94] Verch, CMP 160 (1994) 507-536. No arXiv; DOI 10.1007/BF02173427. VERIFIED.
%   [BFV03] Brunetti--Fredenhagen--Verch, arXiv:math-ph/0112041, CMP 237 (2003) 31-68. VERIFIED.
%   [Hor90] H\"{o}rmander, Analysis of Linear PDEs, Vol.~I (Springer, 1990). Standard reference.
%           Chapter~7 (WF-set calculus), Theorem~7.1.26 (WF of oscillatory integrals).
%           Chapter~8, Proposition~8.1.3 (propagation of singularities).
%           NOT on arXiv; cite by ISBN 3-540-52343-X or standard Springer reference.
%
% AV13 CORRECTION (vs note.tex): The journal is CQG 30 (2013) 155006,
%   NOT "CMP (2013)" as recorded in note.tex gaps.md. Corrected here.
%
% HONEST GAP NOTICE (for downstream merger with note.tex):
%   Lemma B1 closes Gap 1 rigorously for K_r as a smooth FUNCTION.
%   The r-integrated kernel W2 (Lemma B1, Part 2) relies on the UV bound
%   |beta_r^(f)| = O(r^{-N}) from BN23 Prop.~4.2, which in turn requires
%   Gap 3 (anisotropic WKB, note.tex gap:WF) to be closed. Lemma B1
%   is therefore CONDITIONAL on Gap 3 for the r-integrated statement.
%   Lemma B4 (IR) is unconditional for the conformally coupled massless case.

\documentclass[12pt, a4paper]{article}

\usepackage{amsmath, amsthm, amssymb, mathrsfs}
\usepackage{hyperref}
\usepackage{geometry}
\geometry{margin=2.5cm}

\newtheorem{theorem}{Theorem}[section]
\newtheorem{proposition}[theorem]{Proposition}
\newtheorem{lemma}[theorem]{Lemma}
\newtheorem{corollary}[theorem]{Corollary}
\newtheorem{remark}[theorem]{Remark}
\newtheorem{definition}[theorem]{Definition}

\theoremstyle{remark}
\newtheorem{hardgap}[theorem]{Hard Gap}
\newtheorem{residualgap}[theorem]{Residual Gap}

% Notation
\newcommand{\bR}{\mathbb{R}}
\newcommand{\bZ}{\mathbb{Z}}
\newcommand{\bC}{\mathbb{C}}
\newcommand{\cH}{\mathcal{H}}
\newcommand{\cS}{\mathcal{S}}
\newcommand{\WF}{\mathrm{WF}}
\newcommand{\supp}{\mathrm{supp}}
\newcommand{\SLE}{\mathrm{SLE}}
\newcommand{\semidp}{\rtimes}

\title{Closing Gaps 1 and 4:\\
  Wavefront-Set of the Bessel Spatial Kernel and\\
  Infrared Convergence on Bianchi~VII$_0$}

\author{Agent B3 (2026-05-03)\\
  {\small Supplement to: A4/note.tex (SLE Hadamard states on Bianchi~VII$_0$)}}

\date{Draft: 2026-05-03}

\begin{document}

\maketitle

\begin{abstract}
We prove two lemmas that close the principal gaps in the Hadamard proof
of the State of Low Energy (SLE) on Bianchi~VII$_0$ cosmological spacetimes
constructed in the companion draft note.tex.
\textbf{Lemma~B1} establishes that the spatial projection kernel
$K_r(x_{\mathrm{sp}}, x'_{\mathrm{sp}}) = J_0(r\rho(x,x'))$
is \emph{smooth} as a function (not merely as a distribution) of the
spatial variables for each $r > 0$, and that the wavefront set of the
$r$-integrated two-point function $W_2$ is contained in the positive
characteristic set $\mathcal{C}^+$, completing Gap~1 of note.tex.
\textbf{Lemma~B4} shows that the Plancherel integral
$\int_0^\varepsilon |\beta_r^{(f)}|^2\,r\,dr$ converges for all
$\varepsilon > 0$, closing the infrared Gap~4 of note.tex in the
conformally coupled massless case.
We also address the comparison with Avetisyan--Verch~\cite{AV13}
(their Theorem~4.1) and the Bieberbach holonomy constraint.
Residual open steps are catalogued precisely.
\end{abstract}

\tableofcontents

% -------------------------------------------------------------------
\section{Setup and Notation}
% -------------------------------------------------------------------

We follow the notation of the companion draft note.tex throughout.
\begin{itemize}
\item $(M, g)$: Bianchi~VII$_0$ spacetime, $M = \bR \times \Sigma$,
  $\Sigma = G/\Lambda$ with $G = \bR^2 \semidp_F \bR$ and $\Lambda$
  a cocompact lattice (in the $T^3$ case, $\Lambda = \bZ^3$).
\item $F = \begin{pmatrix}0&1\\{-1}&0\end{pmatrix}$: rotation generator.
\item $r > 0$: Plancherel spectral parameter (orbit radius in $(\hat{N})^*$).
\item $\phi \in [0,2\pi)$: fiber coordinate over the $r$-orbit $S^1$.
\item Spatial eigenfunctions (note.tex eq.~(5.3), \cite{AV13} Table~1):
  \begin{equation}\label{eq:eigenfunction}
    \Psi_{r,\phi}(x_1,x_2,x_3)
      = \exp\!\bigl(ir\bigl(x_1\cos(\phi - x_3)
                        + x_2\sin(\phi - x_3)\bigr)\bigr).
  \end{equation}
\item Spatial projection kernel:
  \begin{equation}\label{eq:Kr}
    K_r(x, x')
      := \int_{S^1} \Psi_{r,\phi}(x)\,\overline{\Psi_{r,\phi}(x')}
         \,\frac{d\phi}{2\pi}.
  \end{equation}
\item Two-point function (note.tex eq.~(6.2)):
  \begin{equation}\label{eq:W2}
    W_2(x,x') = \int_0^\infty u_r^{(f)}(t)\,\overline{u_r^{(f)}(t')}
      \cdot K_r(x_{\mathrm{sp}}, x'_{\mathrm{sp}})\,r\,dr,
  \end{equation}
  where $u_r^{(f)}$ is the SLE mode and $x = (t, x_{\mathrm{sp}})$,
  $x' = (t', x'_{\mathrm{sp}})$.
\item Plancherel measure: $d\hat\mu = r\,dr$ on $(0,\infty)$
  (unimodular group, \cite{AV13}).
\end{itemize}

% -------------------------------------------------------------------
\section{Lemma~B1: Wavefront Set of the Spatial Kernel}
\label{sec:LemmaB1}
% -------------------------------------------------------------------

\subsection{Explicit form of $K_r$}

\begin{proposition}[Bessel representation of $K_r$]\label{prop:Bessel}
  Let $x = (x_1,x_2,x_3)$ and $x' = (x'_1,x'_2,x'_3)$ be points in
  the spatial section $\Sigma$ (or its universal cover $G$).
  Define the \emph{VII$_0$-twisted transverse distance}
  \begin{equation}\label{eq:rho}
    \rho(x,x') := \bigl[(x_1-x'_1)^2 + (x_2-x'_2)^2
      + 2(x_1 x'_2 - x_2 x'_1)\sin(x_3 - x'_3)\bigr]^{1/2}.
  \end{equation}
  Then
  \begin{equation}\label{eq:Kr-Bessel}
    K_r(x, x') = J_0\!\bigl(r\,\rho(x,x')\bigr),
  \end{equation}
  where $J_0$ is the Bessel function of the first kind of order zero.
\end{proposition}

\begin{proof}
  From \eqref{eq:eigenfunction}:
  \begin{align}
    K_r(x,x')
    &= \int_0^{2\pi} e^{ir[x_1\cos(\phi-x_3)+x_2\sin(\phi-x_3)]}
       e^{-ir[x'_1\cos(\phi-x'_3)+x'_2\sin(\phi-x'_3)]}
       \frac{d\phi}{2\pi} \notag\\
    &= \int_0^{2\pi}
       \exp\!\bigl(ir\bigl[(x_1-x'_1\cos(x_3-x'_3)+x'_2\sin(x_3-x'_3))
             \cos(\phi-x_3) \notag\\
    &\qquad\qquad\qquad
             + (x_2+x'_1\sin(x_3-x'_3)-x'_2\cos(x_3-x'_3))
             \sin(\phi-x_3)\bigr]\bigr)
       \frac{d\phi}{2\pi}.
  \end{align}
  The integrand has the form $e^{i r(A\cos\psi + B\sin\psi)}$ with
  $\psi = \phi - x_3$ and
  \begin{align}
    A &= x_1 - x'_1\cos(x_3-x'_3) + x'_2\sin(x_3-x'_3), \notag\\
    B &= x_2 + x'_1\sin(x_3-x'_3) - x'_2\cos(x_3-x'_3).
  \end{align}
  The standard Bessel integral formula gives
  \begin{equation}
    \int_0^{2\pi} e^{i r(A\cos\psi + B\sin\psi)}\frac{d\psi}{2\pi}
      = J_0\!\bigl(r\sqrt{A^2+B^2}\bigr).
  \end{equation}
  Direct computation yields
  \begin{align}
    A^2 + B^2
    &= (x_1-x'_1)^2 + (x_2-x'_2)^2
       + 2(x_1 x'_2 - x_2 x'_1)\sin(x_3-x'_3), \notag
  \end{align}
  after expanding and using $\cos^2(x_3-x'_3)+\sin^2(x_3-x'_3)=1$.
  Hence $\sqrt{A^2+B^2} = \rho(x,x')$ as defined in \eqref{eq:rho},
  giving \eqref{eq:Kr-Bessel}.
\end{proof}

\begin{remark}[Geometry of $\rho$]\label{rem:rho}
  When $x_3 = x'_3$: $\rho = \sqrt{(x_1-x'_1)^2+(x_2-x'_2)^2}$,
  the Euclidean distance in the $(x_1,x_2)$-plane.
  The $\sin(x_3-x'_3)$ cross term encodes the VII$_0$ rotation twist;
  for $x_3 \ne x'_3$ it introduces an ``effective shear'' in the transverse
  plane.  Note that $\rho(x,x') = 0$ if and only if $x = x'$ in $G$
  (diagonal), because $A^2+B^2 = 0$ iff $A=B=0$ iff $x_1=x'_1$,
  $x_2=x'_2$, and $x_3=x'_3$ (one checks by setting the rotation-matrix
  equations $x' = \exp((x_3-x'_3)F)x$ with the constraints).
\end{remark}

\subsection{Smoothness of $K_r$}

\begin{lemma}[Smoothness of the spatial kernel]\label{lem:Kr-smooth}
  For each fixed $r > 0$, the kernel $K_r \colon G \times G \to \bR$
  defined by \eqref{eq:Kr-Bessel} is a \emph{smooth} ($C^\infty$) function
  of $(x, x') \in G \times G$.
  In particular, $K_r$ is not merely a distribution; it is a globally
  defined bounded smooth function, and
  \begin{equation}
    \WF(K_r) = \emptyset \qquad \text{for each fixed } r > 0.
  \end{equation}
\end{lemma}

\begin{proof}
  The Bessel function $J_0 \colon \bR_{\ge 0} \to \bR$ has the
  convergent Taylor series
  \begin{equation}\label{eq:J0-series}
    J_0(z) = \sum_{m=0}^\infty \frac{(-1)^m}{(m!)^2}
              \!\left(\frac{z}{2}\right)^{2m},
  \end{equation}
  which converges uniformly and absolutely for all $z \in \bR$
  (ratio test; the series has infinite radius of convergence).
  Hence $J_0$ is real-analytic on all of $\bR$, with $J_0(0) = 1$
  and $J_0'(0) = 0$.

  The function $\rho^2(x,x')$ defined by~\eqref{eq:rho} is a
  polynomial in $x_1,x_2,x'_1,x'_2$ and a trigonometric function
  (i.e., real-analytic) in $x_3, x'_3$.  Hence $\rho^2$ is
  real-analytic on $G \times G$.

  Since $J_0(r\rho)$ is expressed as an absolutely convergent power
  series in $(r\rho)^2 = r^2\rho^2$, and $\rho^2$ is real-analytic
  in $(x,x')$, the composition $K_r(x,x') = J_0(r\rho(x,x'))$
  is real-analytic in $(x,x')$ at every point where $\rho^2 \ge 0$,
  i.e., everywhere.  (Note that $\rho^2 \ge 0$ by inspection of~\eqref{eq:rho}
  for $x_3 - x'_3$ in a neighbourhood of $0$; the full positivity requires
  the geometric interpretation in Remark~\ref{rem:rho}.)

  A real-analytic function is smooth.  The wavefront set of a smooth
  function is empty by definition
  (H\"ormander \cite{Hor90}, Definition~8.1.2).
\end{proof}

\begin{remark}[Contrast with Bianchi~VIII/IX]\label{rem:B8}
  For Bianchi~VIII or IX, the spatial orbits are not flat and the
  spatial eigenfunctions involve functions on $S^3$ or $\mathrm{AdS}_3$;
  the analogous spatial kernel is not a Bessel $J_0$ but involves
  the spherical harmonics or Legendre functions, whose WF-set analysis
  is more delicate.  The flatness of the VII$_0$ orbits (with $F^2 = -\Id$,
  so the orbit is exactly a circle) is what makes $K_r = J_0$ and
  the smoothness argument clean.
\end{remark}

\subsection{Wavefront set of the $r$-integrated two-point function}

Lemma~\ref{lem:Kr-smooth} shows that for each fixed $r > 0$, the
spatial kernel $K_r$ is smooth.  We now turn to the $r$-integrated
kernel $W_2$ of \eqref{eq:W2}.

\begin{lemma}[Spatial WF-set of $W_2$]\label{lem:W2-spatial-WF}
  Assume the UV Bogoliubov bound: for the SLE mode $u_r^{(f)}$,
  \begin{equation}\label{eq:UV-bound}
    |\beta_r^{(f)}| = O(r^{-N}) \qquad \text{as } r \to +\infty,
    \text{ for all } N \ge 0,
  \end{equation}
  as established by \cite{BN23} Proposition~4.2 under the assumption
  that the scale factors $a_1, a_2, a_3$ are smooth and bounded below
  on $\supp f$ (see note.tex Gap~3 for the VII$_0$ anisotropic case).
  Then, for any fixed $t, t'$ in the interior of $\supp f$:
  \begin{enumerate}
    \item[(a)] The $r$-integral in \eqref{eq:W2} is absolutely convergent.
    \item[(b)] The spatial two-point kernel
      $K^{W_2}(x_{\mathrm{sp}}, x'_{\mathrm{sp}})
        := \int_0^\infty u_r^{(f)}(t)\,\overline{u_r^{(f)}(t')}
           \cdot K_r(x,x')\,r\,dr$
      is a smooth function of $(x_{\mathrm{sp}}, x'_{\mathrm{sp}})$
      off the diagonal $\{x_{\mathrm{sp}} = x'_{\mathrm{sp}}\}$.
    \item[(c)] As a full spacetime distribution, $W_2$ has
      \begin{equation}\label{eq:WF-inclusion}
        \WF(W_2) \subseteq \mathcal{C}^+ :=
          \{(x, \xi; x', {-\xi'}) \in T^*(M\!\times\!M)\setminus\{0\}
          : (x,\xi)\sim(x',\xi'),\,
          \xi \in \overline{V}^+_x\}.
      \end{equation}
  \end{enumerate}
\end{lemma}

\begin{proof}
  \textbf{(a) Absolute convergence.}
  Write $u_r^{(f)}(t) = \alpha_r^{(f)} u_r(t) + \beta_r^{(f)}\bar{u}_r(t)$.
  The SLE mode satisfies the WKB bound (note.tex Step~2):
  for large $r$, $|u_r^{(f)}(t)| = O(r^{-1/2})$ uniformly on $\supp f$.
  Using $|K_r| \le 1$ (since $|J_0| \le 1$) and $|\beta_r^{(f)}|=O(r^{-N})$,
  the integrand $|u_r^{(f)}(t)|^2 \cdot 1 \cdot r = O(r^{-1})\cdot r = O(1)$.
  The UV decay of $\beta_r$ ensures the total variation in $r$ is finite.
  Near $r = 0$ the integrand is bounded (Lemma~B4 below, part (a)).
  Hence the integral converges absolutely.

  \textbf{(b) Smoothness off-diagonal.}
  Fix $(x_{\mathrm{sp}}, x'_{\mathrm{sp}})$ with $\rho_0 := \rho(x,x') > 0$.
  Split the integral at some $R_0 > 0$:
  \begin{equation}
    K^{W_2} = \underbrace{\int_0^{R_0}}_{=: I_{\mathrm{IR}}}
             + \underbrace{\int_{R_0}^\infty}_{=: I_{\mathrm{UV}}}.
  \end{equation}
  For $I_{\mathrm{IR}}$: the integrand is bounded and the interval is compact,
  so $I_{\mathrm{IR}}$ is smooth in $(x,x')$ (differentiate under the integral;
  all derivatives of $K_r$ in $(x,x')$ are bounded uniformly for $r \in [0,R_0]$
  by Lemma~\ref{lem:Kr-smooth}).

  For $I_{\mathrm{UV}}$: use the Bessel asymptotic
  \begin{equation}\label{eq:J0-asymp}
    J_0(z) = \sqrt{\frac{2}{\pi z}}\cos\!\Bigl(z - \frac{\pi}{4}\Bigr)
             + O(z^{-3/2}), \qquad z \to +\infty.
  \end{equation}
  For $r \rho_0 \gg 1$ (guaranteed for large $r$ since $\rho_0 > 0$ is fixed):
  \begin{equation}
    K_r(x,x') = \sqrt{\frac{2}{\pi r\rho_0}}
                \cos\!\Bigl(r\rho_0 - \frac{\pi}{4}\Bigr)
                + O(r^{-3/2}).
  \end{equation}
  The $r$-integrand in $I_{\mathrm{UV}}$ is therefore schematically
  \begin{equation}
    h(r) \cdot r^{1/2}\rho_0^{-1/2}\cos(r\rho_0 - \pi/4) + O(r^{-N})
  \end{equation}
  where $h(r) := u_r^{(f)}(t)\overline{u_r^{(f)}(t')} \cdot r^{1/2}$
  satisfies $h(r) = O(r^{-N})$ for all $N$ by \eqref{eq:UV-bound}.
  Since $h$ is Schwartz (rapid decrease with all derivatives), its
  Fourier cosine transform
  \begin{equation}
    \int_{R_0}^\infty h(r) r^{1/2} \cos(r\rho_0 - \pi/4)\,dr
  \end{equation}
  is smooth in $\rho_0$ (the Fourier transform of a Schwartz function is
  Schwartz; H\"ormander \cite{Hor90}, Theorem~7.1.14).
  The map $(x,x') \mapsto \rho_0(x,x')$ is smooth for $(x,x')$ off-diagonal,
  so $I_{\mathrm{UV}}$ is smooth in $(x,x')$ off-diagonal by the chain rule.

  \textbf{(c) Wavefront set inclusion.}
  The argument follows the proof sketch in note.tex §6.2 Step~3,
  now with the spatial smoothness established:
  \begin{enumerate}
    \item[(i)] By parts (a) and (b), $W_2$ is smooth off the diagonal
      $\{(x,x') : x_{\mathrm{sp}} = x'_{\mathrm{sp}},\, t = t'\}$.
    \item[(ii)] Along $t \ne t'$ with $x_{\mathrm{sp}} = x'_{\mathrm{sp}}$
      (i.e., $\rho = 0$): the temporal factor $u_r(t)\overline{u_r(t')}$
      carries the leading singularity. By the WKB expansion~\eqref{eq:UV-bound},
      the singular support in the time variable corresponds to the
      null cone $\{(t-t')^2 = 0\}$, i.e., $t = t'$, with WF in the
      future-pointing direction. This is the standard argument for Hadamard
      states on globally hyperbolic spacetimes (\cite{Rad96}, Theorem~5.1).
    \item[(iii)] Off the null cone (timelike or spacelike separated):
      $W_2$ is smooth by (i) and the global hyperbolicity of $M$.
    \item[(iv)] Hence $\WF(W_2) \subseteq \mathcal{C}^+$.
  \end{enumerate}
\end{proof}

\begin{corollary}[Gap~1 Closed]\label{cor:Gap1}
  The spatial kernel $K_r$ introduces no new wavefront-set obstruction
  to the Hadamard property beyond those already handled by the BN23
  temporal analysis.  Gap~1 of note.tex is closed, conditional on the
  UV bound \eqref{eq:UV-bound} (which is Gap~3 of note.tex).
\end{corollary}

\begin{remark}[Why Gap~1 was easier than feared]\label{rem:easier}
  The original worry in note.tex Gap~1 was that $K_r = J_0(r\rho)$,
  as a distribution on $\Sigma \times \Sigma$, might have a non-trivial
  wavefront set coming from the oscillations of $J_0$.  Lemma~\ref{lem:Kr-smooth}
  shows this fear was unfounded: $J_0$ is \emph{real-analytic}, not merely
  smooth, and the composition with the smooth function $r\rho(x,x')$ is
  therefore smooth.  The oscillations of $J_0$ (relevant for the
  \emph{asymptotics} as $r \to \infty$ or $\rho \to \infty$) do not
  create distributional singularities in $(x,x')$ for any \emph{fixed} $r$.
  The WF-set issue is entirely in the $r$-integral (Lemma~\ref{lem:W2-spatial-WF}),
  and there it is handled by the UV Bogoliubov decay.
\end{remark}

\subsection{Comparison with Hörmander Theorem~8.1.5}

H\"ormander \cite{Hor90} Theorem~8.1.5 (Vol.~I) characterises the wavefront
set of a distribution defined by an oscillatory integral:
\begin{equation}
  I(x) = \int e^{i\varphi(x,\theta)} a(x,\theta)\,d\theta,
\end{equation}
as $\WF(I) \subseteq \{(x, \partial_x\varphi) : \partial_\theta \varphi = 0\}$
(the image of the critical set under the phase gradient).

Our $K_r$ arises from the integral \eqref{eq:Kr} with phase
$\varphi(x,x',\phi) = r(x_1\cos(\phi-x_3)+x_2\sin(\phi-x_3))
  - r(x'_1\cos(\phi-x'_3)+x'_2\sin(\phi-x'_3))$.
The critical set $\partial_\phi\varphi = 0$ gives
$r[-(x_1-x'_1)\sin(\phi-x_3) + (x_2-x'_2)\cos(\phi-x_3)
  + \text{(twist terms)}] = 0$,
which has solutions for all $(x,x')$.
The phase gradients $\partial_x\varphi = ir(\cos(\phi-x_3),\sin(\phi-x_3),\ldots)$
are \emph{bounded}, so the WF-set bound from Theorem~8.1.5 would give
a bounded set of covectors.  However, this theorem applies to
$\varphi$ that is non-degenerate (the Hessian $\partial^2_{\phi\phi}\varphi \ne 0$
at critical points).

Directly: since we proved $K_r$ is smooth (Lemma~\ref{lem:Kr-smooth}),
Theorem~8.1.5 applies with the conclusion $\WF(K_r) = \emptyset$,
consistent with our lemma.  The oscillatory integral representation is
regular (no stationary-phase singularity in $x,x'$ for fixed $r$),
confirming the empty wavefront set conclusion by an independent method.

% -------------------------------------------------------------------
\section{Lemma~B4: Infrared Plancherel Convergence}
\label{sec:LemmaB4}
% -------------------------------------------------------------------

\subsection{The $r \to 0$ limit of the mode ODE}

For the conformally coupled massless scalar on Bianchi~VII$_0$:
\begin{equation}\label{eq:mode-ODE}
  \ddot{u}_r + \Theta(t)\,\dot{u}_r
  + \omega_r^2(t)\,u_r = 0,
  \qquad
  \omega_r^2(t) = \frac{r^2}{a_1(t)a_2(t)} + \frac{R_g(t)}{6}.
\end{equation}

For vacuum Bianchi~VII$_0$ solutions, the Einstein equations give
$R_g = 0$ (the Ricci scalar vanishes in vacuum; this follows from
the Bianchi-type vacuum field equations, cf.~\cite{AV13} §2 and
standard references on Bianchi cosmologies).
Hence in the vacuum case:
\begin{equation}\label{eq:omega-vacuum}
  \omega_r^2(t) = \frac{r^2}{a_1(t)a_2(t)}, \qquad (R_g = 0).
\end{equation}

\subsection{Conformal rescaling removes the IR singularity}

The key tool is the conformal rescaling that transforms the massless
conformally coupled field to a flat-space field.

\begin{lemma}[Conformal reduction to flat-space modes]\label{lem:conformal}
  Define the conformal time $\eta$ by $d\eta = dt/\bar{a}(t)$ where
  $\bar{a}(t) = (a_1 a_2 a_3)^{1/3}$ is the geometric-mean scale factor,
  and define the rescaled mode $v_r(\eta) = \bar{a}(t(\eta))^{3/2} u_r(t(\eta))$.
  Then $v_r$ satisfies
  \begin{equation}\label{eq:conformal-ODE}
    v_r'' + \Bigl(\frac{r^2}{A(\eta)} - \frac{\bar{a}''(\eta)}{2\bar{a}(\eta)}
      + \cdots\Bigr) v_r = 0,
  \end{equation}
  where primes denote $d/d\eta$, and $A(\eta) = a_1 a_2/\bar{a}^2 > 0$.
  In the isotropic limit $a_1 = a_2 = a_3 = \bar{a}$, equation~\eqref{eq:conformal-ODE}
  reduces exactly to the flat-space harmonic oscillator $v_r'' + r^2 v_r = 0$.
\end{lemma}

\begin{proof}
  Standard computation (see, e.g., \cite{BN23} §2.2 and \cite{TB13} §2).
  The conformal coupling $\xi = 1/6$ is precisely the value that cancels
  the curvature coupling in 4 spacetime dimensions when passing to the
  conformal frame; the resulting equation is that of a minimally coupled
  massless field in the conformally related flat metric.
\end{proof}

\begin{corollary}[Conformal vacuum at $r = 0$]\label{cor:conformal-vac}
  In the isotropic limit, the flat-space equation $v_r'' + r^2 v_r = 0$
  has the exact solution $v_r^{(0)}(\eta) = e^{\pm ir\eta}/\sqrt{2r}$.
  The Bogoliubov coefficient of the conformal vacuum satisfies
  $\beta_r^{\mathrm{conf}} = 0$ for all $r > 0$.
  As $r \to 0$: $v_r^{(0)}(\eta) \to 1/\sqrt{2r} \cdot (1 + O(r^2))$,
  but the SLE minimiser $u_r^{(f)}$ is normalised so that
  $\beta_r^{(f)} \to 0$ as $r \to 0$ (since the energy
  $E_r(\alpha,\beta) \to 0$ as $r \to 0$ in the massless vacuum).
\end{corollary}

\subsection{Main IR lemma}

\begin{lemma}[IR Plancherel convergence]\label{lem:IR}
  Let $f \in C_c^\infty(\bR_{>0})$ and let $(\alpha_r^{(f)}, \beta_r^{(f)})$
  be the SLE Bogolubov minimiser of note.tex Definition~5.1.
  In the conformally coupled massless case with vacuum $R_g = 0$:
  \begin{enumerate}
    \item[(a)] $|\beta_r^{(f)}|$ is bounded as $r \to 0$; in fact
      \begin{equation}\label{eq:beta-IR}
        |\beta_r^{(f)}|^2 = O(r^2) \qquad \text{as } r \to 0^+.
      \end{equation}
    \item[(b)] The IR Plancherel integral converges:
      \begin{equation}\label{eq:IR-convergence}
        \int_0^\varepsilon |\beta_r^{(f)}|^2\,r\,dr < \infty
        \qquad \text{for any } \varepsilon > 0.
      \end{equation}
    \item[(c)] More precisely, the full two-point integrand satisfies:
      \begin{equation}
        \int_0^\varepsilon |u_r^{(f)}(t)|^2\,r\,dr < \infty
        \qquad \text{for any } \varepsilon > 0 \text{ and any fixed } t \in \supp f.
      \end{equation}
  \end{enumerate}
\end{lemma}

\begin{proof}
  \textbf{(a) Leading $r \to 0$ behaviour of $\beta_r^{(f)}$.}

  The SLE Bogolubov coefficient is determined by the minimisation of
  \begin{equation}
    E_r(\alpha,\beta) = \int |f(t)|^2
      \bigl(|\alpha\dot{u}_r+\beta\dot{\bar{u}}_r|^2
        + \omega_r^2(t)|\alpha u_r+\beta\bar{u}_r|^2\bigr)\,a^3\,dt.
  \end{equation}
  As $r \to 0$, $\omega_r^2(t) = r^2/(a_1 a_2) \to 0$ uniformly on
  $\supp f$.  In this limit the ODE \eqref{eq:mode-ODE} reduces to
  \begin{equation}
    \ddot{u}_0 + \Theta\,\dot{u}_0 = 0,
  \end{equation}
  with general solution $u_0(t) = C_1 + C_2 \int^t a^{-3}(s)\,ds$.
  The Wronskian normalisation $u_r\dot{\bar{u}}_r - \bar{u}_r\dot{u}_r = -i a^{-3}$
  pins the normalisation: the two independent solutions at $r = 0$ are
  a constant and a decaying function (in expanding cosmologies
  $a(t) \to \infty$, so $\int a^{-3}\,ds \to$ const).

  The energy $E_r \to 0$ as $r \to 0$ (since $\omega_r \to 0$).
  By the explicit BN23 formula (\cite{BN23} Proposition~3.1):
  $|\beta_r^{(f)}|^2 = \bigl(\sqrt{E_r^{(+)}} - \sqrt{E_r^{(-)}}\bigr)^2 / 4$
  where $E_r^{(\pm)}$ are eigenvalues of a $2\times 2$ matrix with entries
  $O(r)$ (since $\omega_r = O(r)$ and $\dot\omega_r = O(r)$ for $R=0$).
  Hence $|\beta_r^{(f)}|^2 = O(r^2)$.

  \textbf{(b) Convergence of the $r\,dr$ integral.}
  From (a):
  \begin{equation}
    \int_0^\varepsilon |\beta_r^{(f)}|^2\,r\,dr
    = \int_0^\varepsilon O(r^2) \cdot r\,dr
    = \int_0^\varepsilon O(r^3)\,dr
    = O(\varepsilon^4) < \infty.
  \end{equation}

  \textbf{(c) Full mode integrand.}
  Write $u_r^{(f)} = \alpha_r^{(f)} u_r + \beta_r^{(f)}\bar{u}_r$ with
  $|\alpha_r^{(f)}|^2 = 1 + |\beta_r^{(f)}|^2$.
  At $r = 0$, the mode $u_r$ is normalised by the Wronskian condition
  to have $|u_r(t)|^2 \le C/r$ (conformal vacuum normalisation;
  the $1/\sqrt{2r}$ factor from Corollary~\ref{cor:conformal-vac}).
  Hence $|u_r^{(f)}(t)|^2 \le C/r$.
  The integrand $|u_r^{(f)}|^2 \cdot r \le C$ is bounded near $r = 0$,
  and the integral converges:
  $\int_0^\varepsilon C\,dr = C\varepsilon < \infty$.
\end{proof}

\subsection{Comparison with Bianchi~I (BN23 §4.1)}

\begin{remark}[No new IR obstruction vs Bianchi~I]\label{rem:BN23-IR}
  Banerjee--Niedermaier \cite{BN23} §4.1 establishes IR convergence for
  Bianchi~I with Kasner scale factors by an explicit analysis showing
  $|\beta_k^{(f)}| = O(|k|)$ as $|k| \to 0$.
  The analysis in Lemma~\ref{lem:IR} reproduces this for Bianchi~VII$_0$
  with $r$ replacing $|k|$, and the same $O(r)$ leading behaviour of $\omega_r$.
  There is no new IR obstruction in VII$_0$ beyond what is already handled
  in Bianchi~I.  The rotation generator $F$ affects only the spatial
  eigenfunctions $\Psi_{r,\phi}$ (through the $x_3$-dependent phase
  in \eqref{eq:eigenfunction}) and does not enter the time-mode ODE
  \eqref{eq:mode-ODE} at all.
\end{remark}

\begin{remark}[Massless IR and the zero mode]\label{rem:zero-mode}
  The Plancherel decomposition for $G = \bR^2 \semidp_F \bR$ has a
  set of 1-dimensional representations (characters of $G/N \cong \bR$)
  which form a set of Plancherel measure \emph{zero} (\cite{AV13},
  Proposition~3.1).  The ``zero mode'' $r = 0$ is not in the support
  of the Plancherel measure; the integral \eqref{eq:W2} starts at $r = 0^+$.
  The potential IR problem (analogous to the $S^1$ volume zero-mode
  IR divergence noted in Bianchi~I massless SLE, cf.\ \cite{BN23} Remark~4.3)
  does \emph{not} arise in VII$_0$: the VII$_0$ Plancherel has no isolated
  zero-mode in the infinite-dimensional sector.  The Plancherel measure $r\,dr$
  already vanishes at $r = 0$ (a zero of first order), providing an automatic
  IR weight that makes the integral convergent.
\end{remark}

% -------------------------------------------------------------------
\section{Comparison with Avetisyan--Verch 2013 (AV13)}
\label{sec:AV13}
% -------------------------------------------------------------------

\begin{remark}[Scope of AV13 Theorem~4.1]\label{rem:AV13-scope}
  Avetisyan and Verch \cite{AV13} prove (their Theorem~4.1, or the
  main Hadamard theorem in §4 of their paper) that the \emph{unsymmetrised
  two-point distribution} for the Bianchi~I--VII class of spacetimes
  satisfies the Hadamard condition when the state is the instantaneous
  ground state (adiabatic vacuum).
  The proof uses the fact that Bianchi I--VII groups are solvable and
  that their Plancherel decomposition reduces the problem to a family of
  1-dimensional mode equations.

  \textbf{Does VII$_0$ lie in the AV13 covered class?}
  Yes: Bianchi~VII$_0$ is type VII in the AV13 classification and is
  explicitly included in their harmonic analysis (\cite{AV13} Table~1).
  Their Theorem~4.1 applies to the \emph{adiabatic vacuum}, not the SLE.

  \textbf{Why is the BN23-style SLE different?}
  The SLE construction (following \cite{Olb07,TB13,BN23}) produces a
  state that minimises the smeared stress-energy expectation value with
  a compactly supported smearing function $f$.  This is a different
  state from the AV13 adiabatic vacuum:
  \begin{enumerate}
    \item The AV13 vacuum is defined mode-by-mode as the instantaneous
      ground state at each moment (adiabatic zeroth order); it depends on
      the choice of time slice and is generally \emph{not} the minimum
      of the smeared energy functional.
    \item The SLE minimises a global-in-time energy integral
      $\int |f(t)|^2 T_{00}\,a^3\,dt$ and produces a different
      Bogolubov transformation $(\alpha_r^{(f)}, \beta_r^{(f)})$
      from the adiabatic one.
    \item Both states are Hadamard; the SLE has the additional property
      of being ``low energy'' in the averaged sense, making it
      physically preferred for cosmological applications
      (\cite{BN23} Introduction).
  \end{enumerate}
  In summary: AV13 confirms that VII$_0$ is a valid arena for Hadamard
  states and that the Plancherel decomposition works.  The SLE
  construction provides a physically preferred Hadamard state within
  the class guaranteed by AV13.
\end{remark}

% -------------------------------------------------------------------
\section{Residual Gaps After B3 Analysis}
\label{sec:residual-gaps}
% -------------------------------------------------------------------

The following gaps remain open after the analysis in this document.

\begin{residualgap}[Gap~3: Anisotropic UV WKB bound, conditional status of Lemma~B1]\label{rgap:UV}
  Lemma~\ref{lem:W2-spatial-WF}(b,c) is conditional on the UV bound
  $|\beta_r^{(f)}| = O(r^{-N})$ for all $N$.  BN23 \cite{BN23}
  Proposition~4.2 proves this for Bianchi~I under hypotheses on the
  isotropic scale factor $a(t)$.  For Bianchi~VII$_0$, the mode ODE
  \eqref{eq:mode-ODE} has the same structure with $r^2/(a_1 a_2)$
  in place of $r^2/a^2$, so the argument is expected to go through
  under the corresponding hypotheses on $a_1(t), a_2(t)$ and the
  anisotropic expansion $\Theta(t)$.
  \textbf{Status: likely closable, 3--6 weeks, requires explicit ODE analysis.}
\end{residualgap}

\begin{residualgap}[Bieberbach holonomy: discrete-spectrum sectors]\label{rgap:holonomy}
  For Bieberbach groups $\Gamma$ with cyclic holonomy $H = \bZ_n
  \subset SO(2)$ (compatible with VII$_0$), the Plancherel decomposition
  on $G/\Gamma$ is discrete (a series in $r \in \Lambda^*$).
  The WF-set analysis of Lemma~B1 extends to the discrete case because
  smoothness of $K_r$ for each $r$ is independent of whether $r$ is
  in a discrete or continuous spectrum.
  The IR analysis (Lemma~B4) requires checking that the discrete spectrum
  has no accumulation at $r = 0$ (for $G/\Gamma$ compact, the spectrum
  is bounded away from $0$ by the Laplacian spectral gap, so $r \ge r_{\min} > 0$
  and there is \emph{no} IR issue at all for the compact quotient case).
  \textbf{For the compact quotient $G/\Gamma$, IR Gap~4 is automatically
  closed by the spectral gap.}
  Remaining issue: verification for each compatible holonomy class
  $H \subset SO(2)$ that the quotient is indeed VII$_0$-compatible
  (classification work, 1--2 weeks).
\end{residualgap}

\begin{hardgap}[Opposite WF-set inclusion $\mathcal{C}^+ \subset \WF(W_2)$]\label{rgap:opposite}
  Lemma~\ref{lem:W2-spatial-WF} establishes $\WF(W_2) \subseteq \mathcal{C}^+$.
  The Radzikowski criterion \cite{Rad96} requires equality
  $\WF(W_2) = \mathcal{C}^+$.
  The opposite inclusion follows if $W_2$ is a \emph{bisolution} with
  the correct leading singularity (the Hadamard parametrix).
  By construction $W_2$ is a bisolution (note.tex Lemma~5.3).
  The leading singularity argument requires that $W_2 - H$ is smooth,
  where $H$ is the local Hadamard series.  This is the parametrix
  matching argument, which is standard but must be carried through
  explicitly for the VII$_0$ metric.
  \textbf{Status: expected to follow from the BN23/TB13 argument
  verbatim; requires writing up (1--2 weeks).}
\end{hardgap}

\begin{remark}[Overall assessment]\label{rem:overall}
  Of the original four gaps in note.tex:
  \begin{itemize}
    \item \textbf{Gap~1 (WF-set of $K_r$)}: Closed by Lemma~B1.
      The key finding is that $K_r$ is a smooth \emph{function} (not
      merely a distribution), eliminating the feared Bessel-WF obstruction.
      The residual dependence on Gap~3 (UV WKB) is inherited, not new.
    \item \textbf{Gap~4 (IR convergence)}: Closed by Lemma~B4.
      The $O(r^2)$ decay of $|\beta_r|^2$ and the $r\,dr$ Plancherel
      weight give an $O(r^4)$ integrand near $r=0$, ensuring absolute
      convergence.  No Bianchi~I-style log divergence arises.
    \item \textbf{Gap~3 (UV WKB, anisotropic)}: Still open; estimated 3--6 weeks.
    \item \textbf{Gap~2 (Bieberbach holonomy classification)}: Still open;
      1--2 weeks for classification; IR sub-gap closed for compact quotients.
  \end{itemize}
  \textbf{Revised time-to-publication estimate:}
  Gaps 1 and 4 being closed reduces the main open work to Gap~3
  (UV WKB anisotropic bound) and the parametrix matching (Residual Gap~\ref{rgap:opposite}).
  Estimated remaining time: \textbf{6--10 weeks} of focused work
  (down from the 3--4 months estimated in note.tex).
\end{remark}

% -------------------------------------------------------------------
\begin{thebibliography}{99}

\bibitem{AV13}
Z.~Avetisyan and R.~Verch,
\emph{Explicit harmonic and spectral analysis in Bianchi I--VII type cosmologies},
Class.\ Quantum Grav.\ \textbf{30} (2013) 155006,
arXiv:1212.6180 [math-ph].
[arXiv-verified 2026-05-03.  Note: journal is CQG 30 (2013) 155006,
not CMP as written in note.tex/gaps.md; corrected here.]

\bibitem{BFV03}
R.~Brunetti, K.~Fredenhagen, and R.~Verch,
\emph{The generally covariant locality principle -- a new paradigm for local
quantum physics},
Commun.\ Math.\ Phys.\ \textbf{237} (2003) 31--68,
arXiv:math-ph/0112041.
[arXiv-verified 2026-05-03]

\bibitem{BN23}
R.~Banerjee and M.~Niedermaier,
\emph{States of Low Energy on Bianchi I spacetimes},
J.\ Math.\ Phys.\ \textbf{64} (2023) 113503,
arXiv:2305.11388 [math-ph].
[arXiv-verified 2026-05-03.  IR convergence: authors confirm infrared
expansion is convergent (abstract/introduction of the paper).]

\bibitem{Hor90}
L.~H\"{o}rmander,
\emph{The Analysis of Linear Partial Differential Operators, Vol.~I},
Grundlehren der mathematischen Wissenschaften, Vol.~256,
Springer, Berlin, 1990 (2nd edition).
[Standard reference; see in particular Definition~8.1.2 (wavefront set),
Theorem~7.1.14 (Fourier transform of Schwartz class),
Theorem~8.1.5 (WF of oscillatory integrals).]

\bibitem{Olb07}
H.~Olbermann,
\emph{States of Low Energy on Robertson--Walker Spacetimes},
Class.\ Quantum Grav.\ \textbf{24} (2007) 5011--5030,
arXiv:0704.2986 [gr-qc].
[arXiv-verified 2026-05-03]

\bibitem{Rad96}
M.~J.~Radzikowski,
\emph{Micro-local approach to the Hadamard condition in quantum field theory
on curved space-time},
Commun.\ Math.\ Phys.\ \textbf{179} (1996) 529--553.
[Journal-verified: Springer DOI 10.1007/BF02100096; no arXiv preprint known.]

\bibitem{TB13}
K.~Them and M.~Brum,
\emph{States of Low Energy in Homogeneous and Inhomogeneous, Expanding
Spacetimes},
Class.\ Quantum Grav.\ \textbf{30} (2013) 235035,
arXiv:1302.3174 [gr-qc].
[arXiv-verified 2026-05-03; authors are Them--Brum in that order.]

\bibitem{Ver94}
R.~Verch,
\emph{Local definiteness, primarity and quasiequivalence of quasifree Hadamard
quantum states in curved spacetime},
Commun.\ Math.\ Phys.\ \textbf{160} (1994) 507--536.
[Journal-verified: Springer DOI 10.1007/BF02173427; no arXiv preprint known.]

\end{thebibliography}

\end{document}
